\exercisesection{Exercises for § 2}

\begin{enumerate}
\def\labelenumi{\arabic{enumi})}
\tightlist
\item
  Let K be a commutative ring with unit element, let E be an A-module, and let \(E^{*}\) be its dual. Denote by \(\varphi\) the canonical homomorphism from \(E \otimes E^{*}\) to End(E).
\end{enumerate}

\begin{enumerate}
\def\labelenumi{\alph{enumi})}
\tightlist
\item
  Any element distinct from 1 in \(\operatorname{End}(\mathbf{E})\) of the form
\end{enumerate}

\[
s _ {x, y ^ {*}} = 1 - \varphi (x \otimes y ^ {*}),
\]

with \(x \in E\) and \(y^{*} \in E^{*}\) , is called a pseudo-reflection in E. Such an element s is called a reflection if x and \(y^{*}\) can be chosen so that \(\langle x, y^{*} \rangle = 2\) ; show that we then have \(s^{2} = 1\) and \(s(x) = -x\) .

\begin{enumerate}
\def\labelenumi{\alph{enumi})}
\setcounter{enumi}{1}
\tightlist
\item
  Let \(x \in \mathbf{E}\) , \(y^{*} \in \mathbf{E}^{*}\) be such that \(\langle x, y^{*} \rangle = 1\) . Show that \(\mathbf{E}\) is the direct sum of the submodule \(Kx\) generated by \(x\) and the orthogonal complement \(H\) of \(y^{*}\) . Show that \(Kx\) is free with basis \(x\) , and that \(s = s_{2x,y^*}\) is equal to 1 on \(H\) and to -1 on \(Kx\) .
\end{enumerate}

\begin{enumerate}
\def\labelenumi{\arabic{enumi})}
\setcounter{enumi}{1}
\tightlist
\item
  With the notation of Exerc. 1, show that \(\det(s_{x,y^*}) = 1 - \langle x, y^* \rangle\) if \(E\) is a free \(K\) -module of finite type.
\end{enumerate}

¶ 3) Let V be a complex Hilbert space with basis \(e_1, \ldots, e_l\) . For \(1 \leqslant i \leqslant l\) , let \(s_i\) be a unitary pseudo-reflection, with vector \(e_i\) , such that \(s_i(e_i) = c_i e_i\) , with \(c_i \neq 1\) ; an element of V is invariant under \(s_i\) if and only if it is orthogonal to \(e_i\) . Let W be the subgroup of \(\mathbf{GL}(V)\) generated by the \(s_i\) .

\begin{enumerate}
\def\labelenumi{\alph{enumi})}
\item
  Let \(i\) be an integer \(\geqslant 1\) . Show that every element of \(\bigwedge^{i}\mathrm{V}\) invariant under W is zero. (Argue by induction on \(l\) ; if \(\mathrm{V}'\) is the subspace of V generated by \(e_1, \ldots, e_{l-1}\) , and if \(e\) is a non-zero vector orthogonal to \(\mathrm{V}'\) , any element of \(\bigwedge^{i}\mathrm{V}\) can be written in the form \(a + (b \wedge e)\) , with \(a \in \bigwedge^{i}\mathrm{V}'\) and \(b \in \bigwedge^{i-1}\mathrm{V}'\) ; if \(a + (b \wedge e)\) is invariant under W, \(a\) and \(b\) are invariant under the subgroup \(\mathrm{W}'\) generated by \(s_1, \ldots, s_{l-1}\) , and the induction hypothesis applies.)
\item
  Assume that W is finite. Show, by using a), that, for any endomorphism A of V,
\end{enumerate}

\[
\begin{array}{c} \sum_ {w \in \mathrm{W}} \det (\mathrm{A} - w) = \operatorname{Card} (\mathrm{W}). \det (\mathrm{A}) \\ \sum_ {w \in \mathrm{W}} \det (1 - \mathrm{A} w) = \operatorname{Card} (\mathrm{W}). \end{array}
\]

Deduce that, for any \(A \in \text{End}(V)\) , there exists \(w \in W\) such that Aw has no non-zero fixed point.

\begin{enumerate}
\def\labelenumi{\alph{enumi})}
\setcounter{enumi}{2}
\item
  Let \(\Gamma\) be the graph whose set of vertices is \([1,l]\) , the arrows being the sets \(\{i,j\}\) such that \(e_{i}\) and \(e_{j}\) are not orthogonal. Show that V is a simple W-module if and only if \(\Gamma\) is connected and non-empty.
\item
  Assume that V is a simple W-module. Show that the W-modules \(\bigwedge^{i}V\) ( \(0 \leqslant i \leqslant l\) ) are simple. (Show that there exists an integer j such that the graph \(\Gamma - \{j\}\) is connected. Argue by induction on l, and apply the induction hypothesis to the subspace \(V'\) generated by the \(e_{i}, i \neq j\) .) Show that these modules are pairwise non-isomorphic \(^{6}\) .
\end{enumerate}

